%% II. Literature review and research gap. Input by access.tex.
\section{Literature Review}
\label{sec:litreview}

\subsection{From Group Differences to the Dynamics of a Session}
Early neuroimaging studies of meditation compared groups: experienced meditators show reduced DMN activity during practice and altered DMN connectivity \cite{Brewer2011,Garrison2015}, and a meta-analysis of 78 functional neuroimaging studies maps the regions engaged by different styles of meditation \cite{Fox2016}. Later work turned to the dynamics within a session. Hasenkamp et al.\ \cite{Hasenkamp2012} segmented focused-attention practice into a cycle of mind wandering, awareness of mind wandering, shifting of attention and sustained focus; Lutz et al.\ \cite{Lutz2008} described focused attention and open monitoring as two strategies of attention regulation; and Christoff et al.\ \cite{Christoff2016} framed mind wandering as spontaneous thought shaped by deliberate and automatic constraints. Predictive-processing accounts describe meditation as a progressive flattening of the predictive hierarchy \cite{Laukkonen2021}, and an active-inference model simulates meta-awareness and attentional control during focused attention \cite{SandvedSmith2021}. These accounts are qualitative or generative, and none of them states how precisely its parameters could be estimated from the data that meditation studies collect.

\subsection{The Brain as a Query Interface Framework}
Among the dynamical accounts, the BQI framework of Vu \cite{BQI} is the one that writes the dynamics of a session as an explicit, parameterised model: it links DMN activity, integrated information and a reported QoC through an ODE system whose rates and gains were fitted to meta-analytic data, and proposes to control them with a closed-loop brain--computer interface \cite[Sec.~7, 9 and 10]{BQI}. Its synthesis of the literature reports a pooled 39.4\,\% reduction of the DMN BOLD signal across 16 studies ($d=1.70$) and a strong anticorrelation between DMN activity and QoC ($r=-0.89$) \cite[Sec.~7.3--7.4]{BQI}. In a companion paper \cite{companion} we showed that the QoC equation of BQI has no sustained level and replaced the system by the state-space model M1. BQI is the theoretical foundation of this study; Section~\ref{sec:theory} states its equations and derives M1 from them.

\subsection{EEG Markers of DMN Activity and of Integration}
The DMN was defined from positron emission tomography and fMRI \cite{Raichle2001}, and its EEG counterparts are indirect. Simultaneous EEG--fMRI studies relate spontaneous fluctuations of band power to resting-state networks: alpha power covaries negatively with a fronto-parietal attention network and 17--23~Hz beta power positively with default-mode regions \cite{Laufs2003}, and each resting-state network has its own combination of EEG rhythms, the DMN being associated with alpha and beta power \cite{Mantini2007}. Posterior alpha power is therefore at best an indirect proxy of DMN activity, and it is used as such here. In meditation research, a systematic review associates mindfulness with increased alpha and theta power \cite{Lomas2015}; mind wandering during breath focus is accompanied by more delta and theta and less alpha and beta power \cite{Braboszcz2011}; experienced practitioners report less mind wandering, with EEG correlates of the reported state \cite{Brandmeyer2018}; and the spectral differences between meditation and mind wandering differ between experienced meditators and novices \cite{RodriguezLarios2021}. BQI identifies elevated theta power as the spectral hallmark of expert meditation \cite[App.~B.2]{BQI}. For integration, $\Phi$ itself is intractable for brain-sized systems; the perturbational complexity index compresses the cortical response to transcranial magnetic stimulation with the Lempel--Ziv algorithm \cite{Casali2013}, the Lempel--Ziv complexity of spontaneous EEG falls under propofol anaesthesia \cite{Schartner2015}, and the aperiodic exponent of the power spectrum offers a spectral alternative \cite{Donoghue2020}. BQI adopts the complexity index as its proxy for $\Phi$ and reports higher values in expert meditators \cite[Sec.~6.3]{BQI}.

\subsection{Identifiability and Experimental Design of Dynamical Models}
Whether the parameters of a dynamical model can be determined from noise-free data is the question of structural identifiability \cite{Bellman1970}; in a regular parametric model a point is locally identifiable if and only if the information matrix is non-singular there \cite{Rothenberg1971}. Practical identifiability asks whether finite, noisy data determine the parameters, and the profile likelihood is its standard tool \cite{Raue2009,Wieland2021}. The Cram\'er--Rao bound \cite{Rao1945,Cramer1946} turns the Fisher information into a lower bound on the variance of any unbiased estimator and underlies optimal experimental design in systems biology \cite{Kreutz2009}. For a latent linear-Gaussian process sampled irregularly, the Kalman filter makes the likelihood exact \cite{Kalman1960,Harvey1989,Durbin2012}. When the model is wrong, the observed and the expected information diverge, which is the basis of the misspecification diagnostic used here \cite{White1982}. These tools had not been applied to experience-sampling designs in meditation before the companion paper \cite{companion}, which did so in simulation.

\subsection{Experience Sampling, Preregistration and Replication}
Probe-caught experience sampling is the standard way to measure mind wandering, and the way probes are framed and timed changes the reported rates \cite{Weinstein2018}. A large replication project found that only about a third of published psychological effects replicated with a significant result \cite{OSC2015}, and preregistration separates confirmatory from exploratory analyses by fixing the plan before the data are seen \cite{Nosek2018}. Open experience-sampling EEG datasets \cite{Delorme2024ds,RodriguezLarios2021} now make preregistered, independent tests of a model possible without new data collection.

\subsection{Research Gap}
\label{sec:gap}
Three gaps follow from this review. First, no design analysis of a dynamical model of meditation has been tested against real experience-sampling EEG; the companion analysis is \emph{in silico}. Second, the parameter values of BQI were fitted to meta-analytic summaries of fMRI contrasts and cross-sectional correlations \cite[Sec.~7.2--7.5]{BQI} and have never been estimated from the EEG time series of single sessions, whose EEG proxies of the DMN and of $\Phi$ are themselves uncalibrated. Third, no rate constant of a meditation model has been replicated across laboratories under a preregistered plan. The present study addresses the three gaps with research questions RQ1--RQ4.
